
The rate $\Theta(\epsilon^{-3})$ is therefore a barrier for $\calC_\theta$.
Faster rates are known in related settings. For smooth NC--C minimax problems,
the optimal deterministic first-order complexity of finding an $\epsilon$-GS
point is $\Theta(\epsilon^{-5/2})$~\cite{pan2026optimal}, in line with the
nonconvex--strongly-concave lower bound with condition number
$\Theta(\epsilon^{-1})$~\cite{zhang2021complexity,li2021complexity}. This rate
is attained by multi-loop methods~\cite{pan2026optimal} and by a single-loop
projected extragradient method with dual gradient
averaging~\cite{zhang2026matching}. For composite minimax and constrained
problems, the methods with rates of order $\epsilon^{-5/2}$ are multi-loop:
nested proximal-point or anchored
schemes~\cite{kong2021accelerated,ke2026efficient}, and inexact ALM and
penalty methods with accelerated inner
solvers~\cite{li2021rate,lin2022complexity,kong2023iteration}. This leads to
our second question:
\begin{quote}
  \textbf{(Q2)}
  \emph{Can a single-loop linearized ALM with a fixed penalty parameter break
  the $\Theta(\epsilon^{-3})$ barrier of $\calC_\theta$ and attain
  $\mathcal O(\epsilon^{-5/2})$ complexity for Problem~\eqref{prob:main}?}
\end{quote}

To answer \textbf{(Q2)}, we introduce Fast LRALM (Algorithm~\ref{alg:b2-prime}
in Section~\ref{sec:acceleration}). Our analysis of \textbf{(Q1)} shows that,
with $\rho$ fixed and $s=\Theta(\epsilon^{-1})$, fixing the dual center at zero
does not change the complexity order; the regularized dual objective then has a
fixed center and is strongly concave. Fast LRALM adds Nesterov extrapolation to
the dual step and keeps one explicit primal--dual update per iteration. We
prove that Fast LRALM finds an $\epsilon$-KKT point of
Problem~\eqref{prob:main} in $\mathcal O(\epsilon^{-5/2})$ iterations with a
fixed $\rho$ (Theorem~\ref{thm:b2-prime-complexity}), which answers
\textbf{(Q2)} affirmatively and matches the rate of the multi-loop methods
above. Because the extrapolated dual step is not projected, the analysis must
bound the multipliers uniformly and control the dual tracking error under a
moving primal center; Section~\ref{sec:acceleration} does both with a single
Lyapunov function. Table~\ref{tab:known-results} compares our guarantees and
parameter scales with representative methods.

